\section{Setting, hierarchies, and standard tools}\label{sec:setting}

This section fixes the sparse structure, the Chebyshev notation, the three
sparse hierarchies studied in the paper, and the standard facts that later
proofs use. We record each fact with the degree bounds used later. The
lemmas in \cref{sec:set-hierarchies,sec:set-tools} are standard; the
private-degree-two hierarchy of \cref{sec:set-recourse} is the specific
relaxation analyzed in \cref{sec:recourse,sec:regularity}.

\subsection{Sparse structure}\label{sec:set-sparse}

\begin{definition}[Junction tree]\label{def:junction-tree}
Let $\Tree$ be a finite tree with nodes $b=1,\dots,t$, and let each node
carry a \emph{bag} $B_b\subseteq\{1,\dots,n\}$. The bags form a
\emph{junction tree} if they cover $\{1,\dots,n\}$ and satisfy the
\emph{running-intersection property}: for every coordinate $i$, the nodes
$b$ with $i\in B_b$ induce a connected subtree of $\Tree$. For an edge
$e=bc\in\Edges$, the \emph{separator} is $S_e=B_b\cap B_c$. We write
$v_b=\abs{B_b}$ and call $w=\max_b v_b$ the \emph{width}.
\end{definition}

In \crefrange{sec:setting}{sec:regularity} and in the box parts of
\cref{sec:extensions,sec:certificates}, the width satisfies $w\ge1$ and the
relaxation order satisfies $r\ge1$. The finite-state models of
\cref{sec:extensions} also allow bags with no continuous coordinates, and
models with no continuous coordinates at all; the conventions for those
cases (empty kernel products, order zero) are stated there and override the
standing assumptions $w\ge1$ and $r\ge1$.

An example whose data involve only private variables can be placed within
the standing positive-width convention by adding one unused shared
coordinate in $[-1,1]$ to its bag. This leaves the original optimum
unchanged. A private functional extends to the rectangular hierarchy by
evaluating that shared coordinate at zero, as in
\cref{prop:convexity-needed}. Individual bags with no shared coordinates are
also allowed when the overall width is positive; their shared polynomial
space consists of constants and their kernel product is one.

Throughout, $x\in[-1,1]^n$ and $x_B=(x_i)_{i\in B}$. The basic problem is
\begin{equation}\label{eq:set-problem}
  \fstar=\min_{x\in[-1,1]^n} f(x),\qquad
  f(x)=\sum_{b=1}^t f_b(x_{B_b}),
\end{equation}
where every $f_b$ is a real polynomial in the variables of its bag. The
decomposition of $f$ into bag polynomials is part of the data: constants in
our bounds depend on it, and different decompositions of the same $f$ can
give different constants. Any product of nondegenerate compact intervals
can be mapped to $[-1,1]^n$ by an invertible affine change of each
coordinate (fixed coordinates are first eliminated); this preserves
bags, total degrees, and agreement of separator moments, so we state all
results on $[-1,1]^n$. When a model is stated on another box we say which
affine change is used.

Running intersection has a convenient consequence for polynomial identities.

\begin{lemma}[Decompositions of zero]\label{lem:zero-decomposition}
Let $p_b\in\Rle{x_{B_b}}{k}$ for $b=1,\dots,t$ satisfy $\sum_b p_b=0$ as a
polynomial in $x$. Then there are polynomials
$\phi_e\in\Rle{x_{S_e}}{k}$, one for each edge, and signs fixed by an
orientation of $\Tree$, such that each $p_b$ is the signed sum of the
$\phi_e$ over the edges incident to $b$.
\end{lemma}

\begin{proof}
Induct on $t$; the case $t=1$ is trivial. Let $\ell$ be a leaf with
neighbor $c$ and edge $e=\ell c$. By running intersection, a coordinate of
$B_\ell\setminus S_e$ belongs to no other bag. Hence every monomial of
$p_\ell$ that involves such a coordinate appears in no other $p_b$, so its
coefficient in $\sum_b p_b=0$ is its coefficient in $p_\ell$, which must be
zero. Thus $p_\ell=:\phi_e\in\Rle{x_{S_e}}{k}$. Replace $p_c$ by
$p_c+\phi_e$, delete the leaf, and apply the induction hypothesis to the
remaining bags; they still satisfy running intersection, and coordinates
that occurred only in the deleted leaf no longer occur and can be
dropped.
\end{proof}

\subsection{Chebyshev calculus}\label{sec:chebyshev}

Let $T_k$ and $U_k$ be the Chebyshev polynomials of the first and second
kind, $T_k(\cos\theta)=\cos(k\theta)$ and
$U_{k-1}(\cos\theta)=\sin(k\theta)/\sin\theta$. They satisfy
$\abs{T_k}\le1$ on $[-1,1]$,
\begin{equation}\label{eq:set-cheb-identities}
  T_jT_k=\tfrac12\bigl(T_{j+k}+T_{\abs{j-k}}\bigr),\qquad
  1-T_k(x)^2=(1-x^2)\,U_{k-1}(x)^2\quad(k\ge1).
\end{equation}
Let $\arcs$ be the normalized arcsine probability measure
$d\arcs(x)=dx/(\pi\sqrt{1-x^2})$ on $[-1,1]$. Then
$\int T_jT_k\,d\arcs=0$ for $j\ne k$, $\int T_0^2\,d\arcs=1$, and
$\int T_k^2\,d\arcs=\tfrac12$ for $k\ge1$. For a bag $B$ we write
$\arcs^B$ for the product measure and, for a multi-index
$\alpha\in\N^B$ (here $\N=\{0,1,2,\dots\}$, so zero indices are allowed), $T_\alpha(x)=\prod_{i\in B}T_{\alpha_i}(x_i)$.

Every polynomial $p$ in the variables of $B$ has a unique expansion
$p=\sum_\alpha c_\alpha T_\alpha$. We use two coefficient functionals:
\begin{equation}\label{eq:set-norms}
  \chebnorm{p}=\sum_\alpha\abs{c_\alpha},\qquad
  \Abud(p)=\sum_\alpha\abs{c_\alpha}\sum_{i\in B}\alpha_i^2 .
\end{equation}
Since $\abs{T_\alpha}\le1$ on the box, $\norm{p}_\infty\le\chebnorm{p}$.
The weighted budget $\Abud(p)$ ignores the constant coefficient and is the
quantity controlled by a kernel whose multipliers satisfy
$1-\kmult_k=O(k^2/m^2)$. For the decomposed objective \eqref{eq:set-problem}
we write
\begin{equation}\label{eq:set-budgets}
  C_b=\chebnorm{f_b-c_{b,0}},\qquad \Cf=\sum_b C_b,\qquad
  \Abud=\sum_b\Abud(f_b),
\end{equation}
where $c_{b,0}$ is the constant Chebyshev coefficient of $f_b$. Thus $C_b$
is the Chebyshev norm of the nonconstant part. Both $\Cf$ and $\Abud$ vanish
exactly when every $f_b$ is constant, and $\Abud\le d^2\Cf$ when every $f_b$
has total degree at most $d$. The oscillation
$\osc(f_b)=\max f_b-\min f_b$ over the bag box satisfies
$\osc(f_b)\le2C_b$.

\begin{lemma}[Products]\label{lem:cheb-product}
For polynomials $p,q$ in the variables of a bag,
$\chebnorm{pq}\le\chebnorm{p}\chebnorm{q}$ and
$\Abud(pq)\le\chebnorm{p}\Abud(q)+\Abud(p)\chebnorm{q}$.
\end{lemma}

\begin{proof}
By \eqref{eq:set-cheb-identities}, $T_\alpha T_\beta$ is the average, over
the $2^{\abs{B}}$ sign choices, of $T_{\gamma}$ with
$\gamma_i\in\{\alpha_i+\beta_i,\abs{\alpha_i-\beta_i}\}$. The average has
coefficient norm at most one, and the average of $\sum_i\gamma_i^2$ equals
$\sum_i(\alpha_i^2+\beta_i^2)$. Expanding $pq$ and using the triangle
inequality proves both bounds.
\end{proof}

\subsection{Truncated cones and the sparse box hierarchies}
\label{sec:set-hierarchies}

For a bag $B$ and an integer $k\ge0$, a sum of squares of degree at most
$2k$ means a finite sum of squares of polynomials of degree at most $k$.
We write $\sos_{2k}(B)$ for these polynomials in the variables $x_B$, and
set $\sos_{2k}(B)=\{0\}$ for $k<0$; terms whose degree allowance is
negative are therefore absent, also when counting matrices.

\begin{definition}[Truncated box cones]\label{def:cones}
The \emph{truncated ordinary box quadratic module} of order $r$ on $B$ is
\[
  \Mod_r(B)=\Bigl\{\sigma_0+\sum_{i\in B}(1-x_i^2)\sigma_i:\
  \sigma_0\in\sos_{2r}(B),\ \sigma_i\in\sos_{2r-2}(B)\Bigr\}.
\]
The \emph{truncated full box preordering} of order $r$ on $B$ is
\[
  \Pre_r(B)=\Bigl\{\sum_{I\subseteq B}\sigma_I\prod_{i\in I}(1-x_i^2):\
  \sigma_I\in\sos_{2r-2\abs{I}}(B)\Bigr\},
\]
where terms with $\abs{I}>r$ are absent. Every summand in either cone has
total degree at most $2r$, and $\Mod_r(B)\subseteq\Pre_r(B)$.
\end{definition}

The ordinary module uses one positive semidefinite (PSD) matrix for the
squares and one for each box generator: $v+1$ matrices for a bag of size
$v$. The preordering adds the products of distinct generators, up to $2^v$
matrices. The difference matters both for the size of the relaxation and,
as \cref{sec:ordinary} shows, for the proof technique.

A linear functional $L:\Rle{x_B}{2r}\to\R$ is \emph{$\Mod_r$-positive} if
$L(p)\ge0$ for all $p\in\Mod_r(B)$, and similarly for $\Pre_r$. Equivalently,
its moment matrix and the localizing matrices of the relevant generator
products are positive semidefinite. Such an $L$ is a \emph{truncated
functional}; it need not be the integral against any measure.

\begin{definition}[Sparse box hierarchies]\label{def:sparse-hierarchy}
Let $r\ge1$ with $\deg f_b\le2r$ for all $b$. A \emph{feasible
ordinary-module family} of order $r$ is a family $(L_b)_{b=1}^t$ of linear
functionals $L_b:\Rle{x_{B_b}}{2r}\to\R$ such that
\begin{enumerate}[label=(\roman*)]
\item $L_b(1)=1$ for every $b$;
\item $L_b$ is $\Mod_r(B_b)$-positive for every $b$;
\item \emph{(exact separator consistency)} $L_b(p)=L_c(p)$ for every edge
  $e=bc$ and every $p\in\Rle{x_{S_e}}{2r}$.
\end{enumerate}
A \emph{feasible preordering family} replaces (ii) by $\Pre_r(B_b)$-positivity.
The \emph{moment values} are
\[
  \rhomod_r=\inf\Bigl\{\sum_bL_b(f_b)\Bigr\},\qquad
  \rhopre_r=\inf\Bigl\{\sum_bL_b(f_b)\Bigr\},
\]
the infima over feasible ordinary-module and preordering families. The
\emph{certificate values} are
\begin{equation}\label{eq:sparse-certificate}
  \lammod_r=\sup\Bigl\{\lambda\in\R:\ f-\lambda=\sum_b q_b,\
  q_b\in\Mod_r(B_b)\Bigr\},
\end{equation}
and $\lampre_r$ is defined in the same way with $\Pre_r(B_b)$. The identity
$f-\lambda=\sum_bq_b$ is an identity of polynomials in $x$.
\end{definition}

Consistency is imposed only along tree edges. By running intersection, two
bags $B_b,B_c$ then agree on every polynomial in $x_{B_b\cap B_c}$ of degree
at most $2r$, because every bag on the tree path from $b$ to $c$ contains
$B_b\cap B_c$.

\begin{lemma}[Weak duality]\label{lem:weak-duality}
For every order $r$,
\begin{equation}\label{eq:weak-duality}
  \lammod_r\le\rhomod_r\le\rhopre_r\le\fstar,\qquad
  \lammod_r\le\lampre_r\le\rhopre_r .
\end{equation}
\end{lemma}

\begin{proof}
Evaluation at a minimizer $x^*$ of \eqref{eq:set-problem},
$L_b(p)=p(x^*_{B_b})$, is a feasible preordering family with objective
$\fstar$, so $\rhopre_r\le\fstar$. Every preordering family is an
ordinary-module family, so $\rhomod_r\le\rhopre_r$; the inclusion of cones
gives $\lammod_r\le\lampre_r$. If $f-\lambda=\sum_bq_b$ with
$q_b\in\Mod_r(B_b)$, then $\sum_b(f_b-q_b)-\lambda=0$, where the constant
$-\lambda$ is assigned to bag~$1$. Apply \cref{lem:zero-decomposition} with
$k=2r$: the left side is a signed sum of separator polynomials, on which
adjacent functionals agree. Hence $\sum_bL_b(f_b)-\lambda=\sum_bL_b(q_b)\ge0$
for every feasible family. The same argument applies to $\Pre_r$.
\end{proof}

Two different statements recur in the paper and should not be confused. A
\emph{rounding statement} applies to every feasible family $(L_b)$: it
constructs a probability law $\nu$ on the feasible set with
$\int f\,d\nu\le\sum_bL_b(f_b)+E$. Taking the infimum gives
$\fstar-\rho_r\le E$. The objective value of an arbitrary feasible family is
\emph{not} a lower bound on $\fstar$; the moment value $\rho_r$ and every
certificate value are. A \emph{certificate statement} asserts that
$f-\fstar+E$ lies in $\sum_b\Mod_r(B_b)$ or $\sum_b\Pre_r(B_b)$, with real
coefficients. For the box hierarchies the two are linked by
\cref{lem:slater} below.

\begin{lemma}[Moment Cauchy--Schwarz]\label{lem:moment-cs}
Let $V$ be a finite-dimensional space of polynomials, $G$ a polynomial, and
$L$ a linear functional with $L(Gp^2)\ge0$ for all $p\in V$. Then
$L(Gpq)^2\le L(Gp^2)\,L(Gq^2)$ for all $p,q\in V$. In particular, if
$L(Gp^2)=0$ then $L(Gpq)=0$ for all $q\in V$.
\end{lemma}

\begin{proof}
The bilinear form $(p,q)\mapsto L(Gpq)$ on $V$ is positive semidefinite.
\end{proof}

\begin{lemma}[Chebyshev moments]\label{lem:cheb-moment}
Let $L:\Rle{x_B}{2r}\to\R$ be $\Mod_r(B)$-positive. For every multi-index
$\alpha$ with $\abs{\alpha}\le r$,
\[
  0\le L(T_\alpha^2)\le L(1),\qquad \abs{L(T_\alpha)}\le L(1).
\]
\end{lemma}

\begin{proof}
Order the coordinates of $B$. By \eqref{eq:set-cheb-identities} and
telescoping,
\begin{equation}\label{eq:set-telescope}
  1-T_\alpha(x)^2=\sum_{i\in B}(1-x_i^2)\,U_{\alpha_i-1}(x_i)^2
  \prod_{j<i}T_{\alpha_j}(x_j)^2,
\end{equation}
where terms with $\alpha_i=0$ are zero. The $i$-th term is $(1-x_i^2)$ times
the square of a polynomial of degree at most $\abs{\alpha}-1\le r-1$. Hence
$1-T_\alpha^2\in\Mod_r(B)$, which gives $L(T_\alpha^2)\le L(1)$, and
$T_\alpha^2\in\Mod_r(B)$ gives $L(T_\alpha^2)\ge0$. \Cref{lem:moment-cs}
with $G=1$, $p=1$, $q=T_\alpha$ completes the proof.
\end{proof}

The lemma is used only for $\abs{\alpha}\le r$, half the degree available
to $L$. Bounds on Chebyshev moments of higher total degree require separate
arguments; one such argument, for the preordering, appears in
\cref{lem:half-degree}.

\begin{lemma}[Compactness]\label{lem:compact}
Let $(L_b)$ be a feasible ordinary-module family of order $r$. Then
$\abs{L_b(x^\beta)}\le1$ for every bag and every monomial with
$\abs{\beta}\le2r$. Consequently the feasible sets of both sparse box
hierarchies are nonempty and compact, and $\rhomod_r$ and $\rhopre_r$ are
attained.
\end{lemma}

\begin{proof}
Let $\abs{\alpha}\le r$ and $\alpha_i>0$. With $q=x^{\alpha-e_i}$, the
generator test gives $L_b((1-x_i^2)q^2)\ge0$, that is,
$L_b(x^{2\alpha})\le L_b(x^{2(\alpha-e_i)})$. Iterating down to $\alpha=0$
gives $0\le L_b(x^{2\alpha})\le1$. Write any $\beta$ with $\abs\beta\le2r$
as $\beta=\alpha+\alpha'$ with $\abs{\alpha},\abs{\alpha'}\le r$;
\cref{lem:moment-cs} gives
$L_b(x^\beta)^2\le L_b(x^{2\alpha})L_b(x^{2\alpha'})\le1$. The feasible sets
are defined by finitely many closed conditions on finitely many bounded
coordinates, and they contain point evaluations.
\end{proof}

\begin{lemma}[Strict feasibility and certificates]\label{lem:slater}
For each bag, let $L_b$ be integration against $\arcs^{B_b}$. This family is
feasible for both sparse box hierarchies, and every moment and localizing
matrix in \cref{def:cones} is positive definite. Consequently
$\lammod_r=\rhomod_r$ and $\lampre_r=\rhopre_r$, and both suprema in
\eqref{eq:sparse-certificate} are attained: $f-\rhomod_r\in\sum_b\Mod_r(B_b)$
and $f-\rhopre_r\in\sum_b\Pre_r(B_b)$, each summand of total degree at most
$2r$, with real coefficients.
\end{lemma}

\begin{proof}
Product arcsine measures have the same marginals on every separator, so the
family is feasible. If $q\ne0$ is a polynomial and $I\subseteq B_b$, then
$q^2\prod_{i\in I}(1-x_i^2)$ is nonnegative and not almost everywhere zero
on the open box, where $\arcs^{B_b}$ has a positive density; its integral is
positive. Hence all moment and localizing matrices are positive definite.

Write the moment problem as a finite conic program in the local moment
vectors, with the PSD constraints of \cref{def:cones}, the normalizations,
and the separator equalities. By \cref{lem:compact} its value is finite.
Its conic dual has a free scalar for the normalization of each bag, a
polynomial $\phi_e\in\Rle{x_{S_e}}{2r}$ for each separator equality, and
asks that each $f_b$ minus its scalar minus the signed separator
polynomials lie in the local cone. Summing these local memberships cancels
the $\phi_e$, so dual feasibility implies a decomposition as in
\eqref{eq:sparse-certificate}; conversely, \cref{lem:zero-decomposition}
turns every decomposition $f-\lambda=\sum_bq_b$ into a dual feasible point
with value $\lambda$. The dual value is therefore $\lammod_r$ (respectively
$\lampre_r$). The conic duality theorem for a strictly feasible primal with
finite value~\cite{bental2001-lectures-modern-convex} gives equality of
the values and attainment of the dual.
\end{proof}

Finite SDP dual attainment, as in \cref{lem:slater}, is different from the
existence of exact polynomial separator functions in an infinite-dimensional
marginal problem. The latter can fail even when every finite SDP is well
posed; \cref{sec:sharpness} gives a quadratic example.

\subsection{Univariate positivity and gluing}\label{sec:set-tools}

\begin{lemma}[Interval positivity with degree bounds]\label{lem:interval-sos}
Let $k\ge0$, and let $p\in\R[x]$ have degree at most $2k$ and satisfy
$p\ge0$ on $[-1,1]$.
Then $p=\sigma_0+(1-x^2)\sigma_1$ with $\sigma_0\in\sos_{2k}$ and
$\sigma_1\in\sos_{2k-2}$.
\end{lemma}

\begin{proof}
If $\deg p=2j$ is even, $j\le k$, the Markov--Lukács
theorem~\cite{powers2000-univariate-interval} gives
$p=\sigma_0+(1-x^2)\sigma_1$ with $\sigma_0\in\sos_{2j}\subseteq\sos_{2k}$
and $\sigma_1\in\sos_{2j-2}\subseteq\sos_{2k-2}$; a nonnegative constant is
the case $j=0$. If $\deg p=2j-1$ is odd, $1\le j\le k$, the same theorem
gives $p=(1+x)\tau_+ +(1-x)\tau_-$ with $\tau_\pm\in\sos_{2j-2}$. The
identity $1\pm x=\tfrac12(1\pm x)^2+\tfrac12(1-x^2)$ converts this into
the claimed form, with $\sigma_0=\tfrac12[(1+x)^2\tau_++(1-x)^2\tau_-]\in
\sos_{2j}\subseteq\sos_{2k}$ and
$\sigma_1=\tfrac12(\tau_++\tau_-)\in\sos_{2j-2}\subseteq\sos_{2k-2}$.
\end{proof}

Zeros of $p$ on the interval are allowed. Multiplying univariate
representations in distinct coordinates gives elements of the preordering,
not of the ordinary module, because products
$(1-x_i^2)(1-x_j^2)$ appear. This is why the full preordering is the
natural cone for product kernels (\cref{sec:kernels}) and why the ordinary
module needs a different kernel (\cref{sec:ordinary}).

All the laws we construct live on products of compact intervals and finite
sets; these are standard Borel spaces, so regular conditional
distributions exist~\cite{kallenberg2002-foundations}.

\begin{lemma}[Gluing on a junction tree]\label{lem:gluing}
Let each coordinate $i$ range over a standard Borel space $X_i$, and let
$\nu_b$ be a probability measure on $X_{B_b}=\prod_{i\in B_b}X_i$ for each
bag. Suppose that for every edge $e=bc$ the marginals of $\nu_b$ and $\nu_c$
on $X_{S_e}$ coincide. Then there is a probability measure $\nu$ on
$\prod_{i=1}^nX_i$ whose marginal on $X_{B_b}$ is $\nu_b$ for every $b$. If
each $\nu_b$ is concentrated on a measurable set $F_b\subseteq X_{B_b}$, then
$\nu$ is concentrated on $\{x:x_{B_b}\in F_b\text{ for all }b\}$.
\end{lemma}

\begin{proof}
Root $\Tree$ and list the bags so that every bag follows its parent. Let
$U_k=B_1\cup\dots\cup B_k$. Suppose $\nu^{(k)}$ is a probability measure on
$X_{U_k}$ with marginals $\nu_1,\dots,\nu_k$. Let $b=k+1$ have parent $c$
and separator $S=B_b\cap B_c$. If $i\in B_b\cap B_j$ with $j\le k$, the tree
path from $b$ to $j$ passes through $c$, so $i\in B_c$ by running
intersection; hence $B_b\cap U_k=S$. Disintegrate
$\nu_b(dx_S,dx_{B_b\setminus S})=\nu_{b,S}(dx_S)\,\kappa(x_S,dx_{B_b\setminus S})$,
and define $\nu^{(k+1)}=\nu^{(k)}(dx_{U_k})\,\kappa(x_S,dx_{B_b\setminus S})$.
Its marginal on $X_{U_k}$ is $\nu^{(k)}$. Its marginal on $X_S$ is that of
$\nu_c$, which equals $\nu_{b,S}$ by assumption; hence its marginal on
$X_{B_b}$ is $\nu_b$. The kernel $\kappa$ is determined only
$\nu_{b,S}$-almost everywhere, and any version gives the same marginals; in
particular an empty separator gives a product. For the last claim,
$\nu(\{x_{B_b}\notin F_b\})=\nu_b(X_{B_b}\setminus F_b)=0$ for each of the
finitely many bags.
\end{proof}

Gluing of consistent local measures under running intersection is
classical in sparse polynomial optimization~\cite{lasserre2006-convergent-sdprelaxations-in-polynomial-optimization}
and for finite distributions~\cite{vorobev1962-consistent-families}. The
issue addressed in this paper is how to obtain consistent local
\emph{measures} from truncated functionals of finite order, with
quantitative control of the objective.

\subsection{Recourse models and the private-degree-two hierarchy}
\label{sec:set-recourse}

\Cref{sec:recourse,sec:regularity} study bags that also carry large private
continuous vectors entering the objective as convex quadratics.

\begin{definition}[Recourse model]\label{def:rec-model}
Let $(\Tree,(B_b))$ be a junction tree on the \emph{shared} coordinates
$x\in[-1,1]^n$. Each bag $b$ also carries a \emph{private} vector
$y_b\in\R^{p_b}$, $p_b\ge0$; private vectors of different bags are
distinct variables. Its \emph{fiber} is
\begin{equation}\label{eq:rec-fiber}
  P_b(x_{B_b})=\bigl\{y\in[-1,1]^{p_b}:\ A_by\le a_b+E_bx_{B_b}\bigr\},
\end{equation}
with real $A_b\in\R^{J_b\times p_b}$, $a_b\in\R^{J_b}$, and
$E_b\in\R^{J_b\times v_b}$. The local objective is
\begin{equation}\label{eq:rec-objective}
  f_b(x,y)=c_b(x)+q_b(x)^{\top}y+y^{\top}Q_b(x)\,y
  =\zy^{\top}H_b(x)\,\zy,\qquad
  \zy=\begin{pmatrix}1\\y\end{pmatrix},\quad
  H_b=\begin{pmatrix}c_b&q_b^{\top}/2\\ q_b/2&Q_b\end{pmatrix},
\end{equation}
with polynomial coefficients in $x_{B_b}$ and symmetric $Q_b$. We assume
\emph{private convexity}: $Q_b(x)\succeq0$ for every $x$ in the shared bag
box. When its feasible set is nonempty, the problem is
\[
  \fstar=\min\Bigl\{\sum_bf_b(x_{B_b},y_b):\ x\in[-1,1]^n,\
  y_b\in P_b(x_{B_b})\text{ for all }b\Bigr\}.
\]
The model has a \emph{fixed private domain} if $E_b=0$ for all $b$, and
\emph{complete recourse} if every fiber is nonempty at every shared box
point. The quantitative theorems of \cref{sec:recourse,sec:regularity}
assume nonempty fixed polytopes or complete recourse, respectively.
\end{definition}

The shared variables carry an arbitrary nonconvex polynomial structure. The
private variables are continuous, belong to one bag, and appear convexly;
only the shared coordinates are coupled across bags. Private convexity is
a pointwise assumption, not an SOS condition on $Q_b$; verifying it may
itself be hard in general, and we use it only where an application supplies
it.

For $I\subseteq B_b$ write $w_I(x)=\prod_{i\in I}(1-x_i^2)$. Let
$\ell_{bj}(x,y)=(a_b+E_bx-A_by)_j$ be the $j$-th row of the fiber
constraints, and let $\epsilon_{bj}=1$ if row $j$ of $E_b$ is nonzero and
$\epsilon_{bj}=0$ otherwise.

\begin{definition}[Private-degree-two hierarchy]\label{def:rec-hierarchy}
Let $r\ge1$ with every entry of $H_b$ of degree at most $2r$. A feasible
family of order $r$ consists of linear functionals
\[
  L_b:\ \Rle{x_{B_b}}{2r}\otimes\Rle{y_b}{2}\longrightarrow\R,
\]
defined on polynomials of shared total degree at most $2r$ and private total
degree at most two, such that, for every $I\subseteq B_b$:
\begin{enumerate}[label=(R\arabic*),start=0]
\item $L_b(1)=1$;
\item $L_b\bigl(w_I\,(s_0(x)+s(x)^{\top}y_b)^2\bigr)\ge0$ for all
  polynomials $s_0,s_1,\dots,s_{p_b}$ of degree at most $r-\abs{I}$;
\item $L_b\bigl(w_I\,s(x)^2\,\ell_{bj}(x,y_b)\bigr)\ge0$ for every row $j$
  and every polynomial $s$ of degree at most $r-\abs{I}-\epsilon_{bj}$;
\item $L_b\bigl(w_I\,s(x)^2\,(1-y_{b,k}^2)\bigr)\ge0$ for $k=1,\dots,p_b$
  and every $s$ of degree at most $r-\abs{I}$;
\item $L_b(p)=L_c(p)$ for every edge $e=bc$ and every
  $p\in\Rle{x_{S_e}}{2r}$.
\end{enumerate}
Conditions with a negative degree allowance are absent. The moment value is
$\rhorec_r=\inf\sum_bL_b(f_b)$ over feasible families.
\end{definition}

Every polynomial tested in (R1)--(R3) has shared degree at most $2r$ and
private degree at most two. Evaluation at a feasible point is a feasible
family, so $\rhorec_r\le\fstar$. The affine bounds $1\pm y_{b,k}\ge0$ are
implied, since
$1\pm y=\tfrac12(1\pm y)^2+\tfrac12(1-y^2)$. The quadratic bounds (R3) are
redundant for the original problem but essential for the relaxation:
without them the relaxation can be unbounded below at every sufficiently
high order (\cref{prop:private-bounds-needed}). Shared positivity uses the full shared
box preordering. The rows with $\epsilon_{bj}=1$ have their degree allowance
reduced by one so that the products in (R2) stay within shared degree $2r$.

The point of the hierarchy is its size. The PSD matrix in (R1) indexed by
$I$ has
\begin{equation}\label{eq:rec-size}
  (p_b+1)\binom{r-\abs{I}+v_b}{v_b}
\end{equation}
rows, the localizing matrices in (R2)--(R3) have at most
$\binom{r-\abs{I}+v_b}{v_b}$ rows, there are $2^{v_b}$ choices of $I$, and the
number of local moments is at most
$\binom{p_b+2}{2}\binom{2r+v_b}{v_b}$. For fixed width these sizes are
polynomial in $p_b$, $J_b$, and $r$, and the exponent of $r$ depends only on
the number $v_b$ of shared coordinates in the bag. By contrast, a standard
total-degree hierarchy on all $v_b+p_b$ variables of the bag has a moment
matrix with $\binom{r+v_b+p_b}{r}$ rows. For $v_b=2$, $p_b=50$, and $r=6$
these numbers are $1{,}428$ and $40{,}475{,}358$. The constants in our error
bounds, however, can grow with the private dimension through the
coefficients of $H_b$ and the fiber data; the hierarchy is not
dimension-free in accuracy.

Partial lifting of selected nonlinear variables is an established way to
reduce relaxation size~\cite{kahl2005-globally-optimal-estimates}. Holding
an SOS-convex decision-variable module at fixed degree while the
uncertainty-variable order grows is also used in robust polynomial matrix
inequality hierarchies~\cite{wang2025-a-moment-sum-of-squares}. Here we give
a quantitative analysis of the specified sparse private-degree-two hierarchy.

The private-degree-two hierarchy keeps the private degree at two while
retaining shared degree up to $2r$ in mixed moments, so its moments have
joint degree up to $2r+2$. A total-degree hierarchy of order $r$ retains all
joint moments of degree up to $2r$, including higher private degrees; if it
uses a full preordering, it also contains products of distinct private
constraints. We do not compare the two formulations by inclusion. Lower
bounds that we state for both are proved separately for each, by actual
local measures, which are feasible for every hierarchy considered here.
